% ---------------------------------------------------------------------------
\section{Discussion}
\label{sec:discussion}

\subsection{What the generator shows}

The voynichizer is a constructive argument. Its books reproduce most of the measured writing rules of the manuscript, to
the point that two classifiers retrained on each book separate its pages from the real ones only slightly better than
chance, and every book carries an encrypted text that comes back exactly with the key. Two claims must be kept apart.
\begin{itemize}
\item \textbf{A book with a message and a book without one} are samples of the same distribution, because the message
  enters only as encrypted bits, which are indistinguishable from the random filler that drives a book without a message.
  No test on the text can separate them without the key. This holds by construction, whatever the judge, as long as the
  cipher is sound; the measurement on version~21 is in Section~\ref{sec:evaluation}.
\item \textbf{A generated book and the real manuscript} are not the same distribution. Our judges separate them a
  little (AUC 0.55 and 0.60), several properties are not reproduced, and a stronger or independently built judge would
  probably do better. How close the book comes is a measured fact about our instruments.
\end{itemize}
Together they mean that statistical regularities of the kind measured here, however strong, cannot decide by themselves
whether the Voynich manuscript carries a message: a text written under these rules can carry one, and whether it does
cannot be read from its statistics. This does not show that the manuscript carries a message. It shows that the debate
between ``meaningless'' and ``enciphered'' cannot be settled by the statistics of the text alone
\citep[compare][]{gaskellbowern2022,bowerngaskell2022,greshko2025,kinnison2026}.

\subsection{What the measurements exclude and what they leave open}

\paragraph{Excluded, with the controls described}
\begin{itemize}
\item a natural language written in an ordinary alphabet and read with a simple or homophonic substitution, in the
  languages tested;
\item the published readings we could test;
\item the historical encryption mechanisms we simulated, one at a time;
\item a text invented by hand without rules, as the volunteers' gibberish is;
\item that the rules of the manuscript are the ordinary habits of the 79 Nordic and German manuscripts we measured
  (with the reservations on the short state in Section~\ref{sec:scribes}).
\end{itemize}

\paragraph{Open} Four scenarios, which our data do not separate:
\begin{enumerate}
\item \textbf{A language written with a very artificial writing system}, in which the variant forms carry no meaning but
  follow rules of form and spaces are placed by rules of form. The junction rules are compatible with a phonetically
  written language: two of the 24 languages, Sanskrit and Tagalog, have a stronger junction.
\item \textbf{A code or cipher of an untested kind}, for example with a repertory of words or a lost key book. The
  voynichizer shows one way in which a message can hide in a text with these rules.
\item \textbf{No message}: a text produced by a procedure that can be executed by hand. Generators of this kind
  reproduce many basic properties, but none of those we measured reproduces the junction rules, the closed line and the
  margin avoidance together.
\item \textbf{A line that is a unit of content} (a verse, an entry, a recipe, a formula). It would explain at once the
  closed line, the line-edge forms and the register of first and last lines; a herbal printed one verse per line has a
  closed line too. It is compatible with any of the first three scenarios.
\end{enumerate}

\subsection{What would change our mind}

\begin{itemize}
\item A reading that passes the controls used here: one that works on the real text and not on shuffled text or on
  text generated without a message, that produces meaningful sentences on pages not used to build it, and that explains
  the junction rules, the closed line, the copying from the line above and the margin avoidance.
\item A medieval manuscript, or a historical cipher, with the same rules; in particular Latin or Italian manuscripts of
  the fifteenth century, which the battery lacks.
\item A simple procedure, executable by hand, that reproduces them all together.
\end{itemize}

\noindent A reading will need an external anchor that regularities cannot give: a text of which the manuscript is a copy,
a secure identification of many plants to use as cribs, a key found in an archive.

% ---------------------------------------------------------------------------
\section{Limitations}
\label{sec:limits}

\begin{itemize}
\item \textbf{Preregistration and exploration.} The preregistrations are internal, dated files in a public
  version-controlled history, without third-party timestamps. The programme was exploratory, and confirmations on data
  not used to form a hypothesis are few (Section~\ref{sec:methods}).
\item \textbf{Transliterations.} Everything rests on transliterations that share the line segmentation. Without checking
  the images we cannot separate what the scribe did from what the transcribers read, for example the \eva{q} that drops
  after a drawing; for uncertain spaces, the image measurements of \citet{rozanova2026} show that they are physically
  narrower.
\item \textbf{Comparison manuscripts.} Old Norse and Early New High German only; no fifteenth-century Latin or Italian
  manuscript transcribed at the level of letter forms. The files do not mark changes of hand, and the 73 German
  manuscripts were measured together for the state of the choices, so a single scribe's state could be diluted. For the
  short state the power is insufficient in 13 of 29 choices.
\item \textbf{Comparison languages.} 24 natural languages, with lines made by cutting prose; 10{,}000 words per text.
\item \textbf{Historical ciphers.} Only one real enciphered manuscript (the Copiale); the DECODE archive of early modern
  ciphers requires an account and has no reuse licence, and we found no fifteenth-century cipher there with both
  transcription and key.
\item \textbf{Gaps.} Three tests could not be run because their measures did not separate the cases on synthetic texts:
  whether the state of the choices lasts a number of words or a span of writing (glyphs), whether the states of
  different choices change at the same points, and whether the short state crosses the line break. Whether the state
  survives the gap of a drawing is undecided for lack of data.
\item \textbf{The generator.} The judges and the scorecard are the project's own, and the generator was tuned against
  them; an independently built judge has not been tried. Keys 1--12 were used during development. The books are not
  indistinguishable from the manuscript. The encryption has not been reviewed by a cryptographer, and the scrypt salt is
  fixed.
\item \textbf{Effect sizes.} Several of the properties are small in absolute terms (the short state is about 13
  percentage points of agreement between adjacent words); the look-ahead and the drift of single choices other than
  \eva{sh}/\eva{ch} do not survive the strictest multiplicity correction.
\end{itemize}

% ---------------------------------------------------------------------------
\section{Availability}

\begin{itemize}
\item \textbf{Research record} \citep{repository}: preregistrations (\texttt{preregistrazioni/}), programs
  (\texttt{esperimenti/}), results with provenance records (\texttt{risultati/}), the laboratory notebook
  (\texttt{QUADERNO.md}, in Italian), the log of all experiments (\url{white_paper/revisione/registro_esperimenti.csv})
  and the sources of this paper with the script that draws its figures. Re-running an experiment gives the same numbers.
  Copyright-protected corpora are not included and must be downloaded from their sources (Appendix~\ref{app:sources}).
\item \textbf{The voynichizer} \citep{voynichizer}: program (MIT licence), the transliteration it uses (CC0), the font, and
  a website that runs version~21 in the browser. The example book of Figure~\ref{fig:page} and its message are in the
  research repository (\url{white_paper/en/figure/}).
\end{itemize}
